\section*{SUBALGEBRAS}

\begin{enumerate}
\def\labelenumi{\arabic{enumi})}
\setcounter{enumi}{12}
\item
  If \(P \subset R\) , put \(g^{P} = \bigoplus_{\alpha \in P} g^{\alpha}\) and \(h_{P} = \sum_{\alpha \in P} kH_{\alpha}\) . Let \(P \subset R\) , \(h'\) a vector subspace of h, and \(a = h' \oplus g^{P}\) . Then a is a subalgebra of g if and only if P is a closed subset of R and \(h'\) contains \(h_{P \cap (-P)}\) ; a is reductive in g if and only if P = -P; and a is solvable if and only if \(P \cap (-P) = \varnothing\) .
\item
  Let \(\mathbf{P}\) be a closed subset of \(\mathbf{R}\) , and \(\mathfrak{b} = \mathfrak{h} \oplus \mathfrak{g}^{\mathrm{P}}\) . The following conditions are equivalent:
\end{enumerate}

\begin{enumerate}
\def\labelenumi{(\roman{enumi})}
\item
  \(\mathfrak{b}\) is a maximal solvable subalgebra of \(\mathfrak{g}\) ;
\item
  \(\mathrm{P} \cap (-\mathrm{P}) = \varnothing\) and \(\mathrm{P} \cup (-\mathrm{P}) = \mathrm{R}\) ;
\item
  there exists a chamber \(\mathbf{C}\) of \(\mathbf{R}\) such that \(\mathrm{P} = \mathrm{R}_{+}(\mathrm{C})\) (cf.~Chap. VI, §1, no. 6).
\end{enumerate}

A Borel subalgebra of \((\mathfrak{g},\mathfrak{h})\) is a subalgebra of g containing h and satisfying the above conditions. A subalgebra b of g is called a Borel subalgebra of g if there exists a splitting Cartan subalgebra \(h'\) of g such that b is a Borel subalgebra of \((\mathfrak{g},\mathfrak{h}')\) ; if k is algebraically closed, this is equivalent to saying that b is a maximal solvable subalgebra of g.

Let \(\mathfrak{b} = \mathfrak{h} \oplus \mathfrak{g}^{\mathrm{R}_{+}(C)}\) be a Borel subalgebra of \((\mathfrak{g}, \mathfrak{h})\) . The largest nilpotent ideal of \(\mathfrak{b}\) is \([\mathfrak{b}, \mathfrak{b}] = \mathfrak{g}^{\mathrm{R}_{+}(C)}\) . Let \(B\) be the basis of \(R\) associated to \(C\) ; the algebra \([\mathfrak{b}, \mathfrak{b}]\) is generated by the \(\mathfrak{g}^{\alpha}\) for \(\alpha \in B\) .

If \(\mathfrak{b},\mathfrak{b}'\) are Borel subalgebras of \(\mathfrak{g}\) , there exists a Cartan subalgebra of \(\mathfrak{g}\) contained in \(\mathfrak{b} \cap \mathfrak{b}'\) ; such a subalgebra is splitting.

\begin{enumerate}
\def\labelenumi{\arabic{enumi})}
\setcounter{enumi}{14}
\tightlist
\item
  Let \(\mathbf{P}\) be a closed subset of \(\mathbf{R}\) , and \(\mathfrak{p} = \mathfrak{h} \oplus \mathfrak{g}^{\mathrm{P}}\) . The following conditions are equivalent:
\end{enumerate}

\begin{enumerate}
\def\labelenumi{(\roman{enumi})}
\item
  p contains a Borel subalgebra of \((\mathfrak{g}, \mathfrak{h})\) ;
\item
  \(P \cup (-P) = R;\)
\item
  there exists a chamber C of R such that \(\mathrm{P} \supset \mathrm{R}_{+}(\mathrm{C})\) .
\end{enumerate}

A parabolic subalgebra of \((\mathfrak{g},\mathfrak{h})\) is a subalgebra of g containing h and satisfying the above conditions. A subalgebra p of g is called parabolic if there exists a splitting Cartan subalgebra \(h'\) of g such that p is a parabolic subalgebra of \((\mathfrak{g},\mathfrak{h}')\) .

Let \(\mathfrak{p} = \mathfrak{h} \oplus \mathfrak{g}^{\mathrm{P}}\) be a parabolic subalgebra of \((\mathfrak{g}, \mathfrak{h})\) , Q the set of \(\alpha \in \mathrm{P}\) such that \(-\alpha \notin \mathrm{P}\) , and \(\mathfrak{s} = \mathfrak{h} \oplus \mathfrak{g}^{\mathrm{P} \cap (-\mathrm{P})}\) . Then \(\mathfrak{p} = \mathfrak{s} \oplus \mathfrak{g}^{\mathrm{Q}}\) , \(\mathfrak{s}\) is reductive in \(\mathfrak{g}\) , and \(\mathfrak{g}^{\mathrm{Q}}\) is the largest nilpotent ideal of \(\mathfrak{p}\) and the nilpotent radical of \(\mathfrak{p}\) . The centre of \(\mathfrak{p}\) is 0.

